\documentclass{article}
\usepackage{titlesec}

\usepackage[citestyle=authortitle-terse,backref]{biblatex}
\addbibresource{bibliography.bib}

\usepackage[left=4cm, right=4cm]{geometry}
\usepackage{palatino,eulervm,dutchcal,xcolor}
\usepackage{graphicx,subcaption,float}
\usepackage{enumitem,parskip,multicol,array}
\newcolumntype{L}{>{$}c<{$}}%table in mathmode
\usepackage{amsthm,amssymb,amsmath,mathtools,thmtools,tikz,tikz-cd}
\usetikzlibrary{%
	matrix,%
	calc,%
	arrows,%
	shapes,
	decorations.markings,backgrounds,calc,intersections}
\tikzcdset{scale cd/.style={every label/.append style={scale=#1},
		cells={nodes={scale=#1}}}}
\usepackage[bookmarks,bookmarksopen,bookmarksdepth=3]{hyperref}
\hypersetup{%colores
	colorlinks=true,
	urlcolor=blue,
	linkcolor=magenta,
	citecolor=blue,
	filecolor=blue,
	urlbordercolor=white,
	linkbordercolor=white,
	citebordercolor=white,
	filebordercolor=white}
\usepackage{cleveref}

\definecolor{blue-violet}{rgb}{0.54, 0.17, 0.89}
\definecolor{azure}{rgb}{0.0, 0.5, 1.0}
\definecolor{green(ncs)}{rgb}{0.0, 0.62, 0.42}
\definecolor{forestgreen}{rgb}{0.13, 0.55, 0.13}
\definecolor{limegreen}{rgb}{0.2, 0.8, 0.2}
\definecolor{palatinateblue}{rgb}{0.15, 0.23, 0.89}
\definecolor{trueblue}{rgb}{0.0, 0.45, 0.81}
\definecolor{goldenyellow}{rgb}{1.0, 0.87, 0.0}
\definecolor{fashionfuchsia}{rgb}{0.96, 0.0, 0.63}
\definecolor{brightcerulean}{rgb}{0.11, 0.67, 0.84}
\definecolor{jonquil}{rgb}{0.98, 0.85, 0.37}
\definecolor{lavendermagenta}{rgb}{0.93, 0.51, 0.93}
\definecolor{peru}{rgb}{0.8, 0.52, 0.25}
\definecolor{persimmon}{rgb}{0.93, 0.35, 0.0}
\definecolor{persianred}{rgb}{0.8, 0.2, 0.2}
\definecolor{persianblue}{rgb}{0.11, 0.22, 0.73}
\definecolor{persiangreen}{rgb}{0.0, 0.65, 0.58}
\definecolor{persianyellow}{rgb}{0.9, 0.89, 0.0}



\declaretheoremstyle[headfont=\color{trueblue}\normalfont\bfseries,]{colored1}
\declaretheoremstyle[headfont=\color{forestgreen}\normalfont\bfseries,]{colored2}
\declaretheoremstyle[headfont=\color{peru}\normalfont\bfseries,]{colored3}
\declaretheoremstyle[headfont=\color{persiangreen}\normalfont\bfseries,]{colored4}
\declaretheoremstyle[headfont=\color{brightcerulean}\normalfont\bfseries,]{colored5}
\declaretheoremstyle[headfont=\color{lavendermagenta}\normalfont\bfseries,]{colored6}
\declaretheoremstyle[headfont=\color{blue-violet}\normalfont\bfseries,]{colored7}
\declaretheoremstyle[headfont=\color{green(ncs)}\normalfont\bfseries,]{colored8}
\declaretheoremstyle[headfont=\color{peru}\normalfont\bfseries,]{colored9}
\declaretheoremstyle[headfont=\color{persiangreen}\normalfont\bfseries,]{colored10}

\declaretheorem[style=colored1,numbered=no,name=Theorem]{thm}
\declaretheorem[style=colored2,numbered=no,name=Proposition]{prop}
\declaretheorem[style=colored3,numbered=no,name=Lemma]{lemma}
\declaretheorem[style=colored4,numbered=no,name=Corollary]{coro}
\declaretheorem[style=colored5,numbered=no,name=Example]{example}
\declaretheorem[style=colored5,numbered=no,name=Examples]{exemplos}
\declaretheorem[style=colored7,numbered=no,name=Exercise]{exercise}
\declaretheorem[style=colored6,numbered=no,name=Remark]{remark}
\declaretheorem[style=colored8,numbered=no,name=Claim]{claim}
\declaretheorem[style=colored9,numbered=no,name=Definition]{defn}
\declaretheorem[style=colored10,numbered=no,name=Question]{question}

\numberwithin{equation}{section}

\newcommand{\A}{\mathbb{A}}
\newcommand{\R}{\mathbb{R}}
\newcommand{\Z}{\mathbb{Z}}
\newcommand{\N}{\mathbb{N}}
\newcommand{\C}{\mathbb{C}}
\newcommand{\Q}{\mathbb{Q}}
\newcommand{\D}{\mathbb{D}}
\renewcommand{\P}{\mathbb{P}}

\newcommand{\Ac}{\mathcal{A}}
\newcommand{\Bc}{\mathcal{B}}
\newcommand{\Cc}{\mathcal{C}}
\newcommand{\Dc}{\mathcal{D}}
\newcommand{\Ec}{\mathcal{E}}
\newcommand{\Fc}{\mathcal{F}}
\newcommand{\Gc}{\mathcal{G}}
\newcommand{\Hc}{\mathcal{H}}
\newcommand{\Ic}{\mathcal{I}}
\newcommand{\Jc}{\mathcal{J}}
\newcommand{\Kc}{\mathcal{K}}
\newcommand{\Lc}{\mathcal{L}}
\newcommand{\Mc}{\mathcal{M}}
\newcommand{\Nc}{\mathcal{N}}
\newcommand{\Oc}{\mathcal{O}}
\newcommand{\Pc}{\mathcal{P}}
\newcommand{\Qc}{\mathcal{Q}}
\newcommand{\Rc}{\mathcal{R}}
\newcommand{\Sc}{\mathcal{S}}
\newcommand{\Tc}{\mathcal{T}}
\newcommand{\Uc}{\mathcal{U}}
\newcommand{\Vc}{\mathcal{V}}
\newcommand{\Wc}{\mathcal{W}}
\newcommand{\Xc}{\mathcal{X}}
\newcommand{\Yc}{\mathcal{Y}}
\newcommand{\Zc}{\mathcal{Z}}

\newcommand{\mf}{\mathfrak{m}}

\DeclareMathOperator{\img}{img}
\DeclareMathOperator{\Arg}{Arg}
\DeclareMathOperator{\id}{id}
\DeclareMathOperator{\pt}{pt}
\DeclareMathOperator{\Alt}{Alt}
\DeclareMathOperator{\sgn}{sgn}
\DeclareMathOperator{\hTop}{h-Top}
\DeclareMathOperator{\supp}{supp}
\DeclareMathOperator{\Int}{Int}
\DeclareMathOperator{\sing}{sing}
\DeclareMathOperator{\Ob}{Ob}
\DeclareMathOperator{\Mor}{Mor}
\DeclareMathOperator{\Top}{Top}
\DeclareMathOperator{\Set}{Set}
\DeclareMathOperator{\CGWH}{CGWH}
\DeclareMathOperator{\Hom}{Hom}
\DeclareMathOperator{\Map}{Map}
\DeclareMathOperator{\Tot}{Tot}
\DeclareMathOperator{\op}{op}
\DeclareMathOperator{\ev}{ev}
\DeclareMathOperator{\hofib}{hofib}
\DeclareMathOperator{\rel}{rel}
\DeclareMathOperator{\Nat}{Nat}
\DeclareMathOperator{\Arr}{Arr}
\DeclareMathOperator{\Funct}{Funct}
\DeclareMathOperator{\colim}{colim}
\DeclareMathOperator{\GL}{GL}
\DeclareMathOperator{\SL}{SL}
\DeclareMathOperator{\SO}{SO}
\DeclareMathOperator{\U}{U}
\DeclareMathOperator{\SU}{SU}
\DeclareMathOperator{\Sp}{Sp}
\DeclareMathOperator{\M}{M}
\DeclareMathOperator{\Aut}{Aut}
\DeclareMathOperator{\PGL}{PGL}
\DeclareMathOperator{\PSL}{PSL}
\DeclareMathOperator{\Vect}{Vect}
\DeclareMathOperator{\VectBund}{VectBund}
\DeclareMathOperator{\Open}{Open}
\DeclareMathOperator{\coker}{coker}

\begin{document}
	{\Large de Rham and Dolbeault theorems}
	\subsection*{Preliminaries}
	\begin{defn}If $\Ac$ is a sheaf of rings, a sheaf $\Fc$ of $\Ac$-modules is said to be a \textbf{\textit{sheaf of free $\Ac$-modules}} if there exists an integer $n$ such that $\Fc$ is locally isomorphic to $\Ac^n$ as a sheaf of $\Ac$-modules. The integer $n$ is called the \textbf{\textit{rank}} of $\Fc$.
	\end{defn}
	\begin{lemma}
		Let $\Ac$ be the sheaf of continuous, differentiable or holomorphic functions over a topological, differentiable or complex manifold $X$. For every vector topological, differentiable or holomorphic bundle $E\to X$, the presheaf $\Ec$
		\[U\mapsto\{\text{cont., dif. or hol. sections }\sigma:U\to E|_U\}\]
		is a sheaf of modules over its corresponding sheaf of functions such that the correspondence
		\[E\mapsto\Ec\]
		is an equivalence of categories between vector bundles and sheaves of free $\Ac$-modules.
	\end{lemma}
	\begin{defn}
			\begin{itemize}
			\item A \textbf{\textit{morphism of sheaves}} $\phi:\Fc\to\Gc$ is a collection, for each open set $U$, of morphisms $\phi_U:\Fc(U)\to\Gc(U)$ such that for $\sigma\in\Fc(U)$ and $V\subset U$ we have $\phi_U(\sigma)|_V=\phi_V(\sigma|_V)$.
			\item The morphism $\phi$ is \textbf{\textit{injective}} (resp. \textbf{\textit{surjective}}) if for every $x\in X$ the morphism
			\[\phi_x:\Fc_x\to\Gc_x\]
			is injective (resp. surjective).
			\item The presheaf $U\mapsto \ker(\phi_U:\Fc(U)\to\Gc(U))$ is a sheaf called the \textbf{\textit{kernel}} of $\phi$.
			\item The sheaf associated to the presheaf $U\mapsto\img(\phi_U:\Fc(U)\to\Gc(U))$ is called \textbf{\textit{image}} of $\phi$.
			\item The sheaf associated to the presheaf $U\mapsto\coker(\phi_U:\Fc(U)\to\Gc(U))$ is called \textbf{\textit{coker}} of $\phi$.
		\end{itemize}
		\end{defn}
	\subsection*{de Rham's theorem}
	
	Let's prove that singular cohomology groups are isomorphic to de Rham cohomology groups:
	\[H_{dR}^q(X)\cong H^q_{\sing}(X,\R).\]
	References are from \cite{voisin}. {\color{blue-violet}Purple} means I'm unsure.
	
	
	\begin{enumerate}
		\item \begin{lemma}[4.26]
			If an abelian category $\Cc$ has sufficiently many injective objects, then every every object of $\Cc$ admits an injective resolution, i.e. a resolution $I^\bullet$ by a complex all of whose objects are injective objects.
		\end{lemma}
		\item 	\begin{thm}[4.28, Derived functors of a left-exact functor exist.]
			Let $\Cc$ and $\Cc'$ be two abelian categories, and let $F$ be a left-exact functor from $\Cc$ to $\Cc'$. Assume that $\Cc$ has sufficiently many injective objects. For every object $M$ of $\Cc$ there exist objects $R^iF(M)$ for $i\geq0$ in $\Cc'$ satisfying the following conditions:
			\begin{itemize}
				\item (Conditions...)
			\end{itemize}
		\end{thm}
		\item \begin{defn}[4.31]
			We say that an object $M$ of $\Cc$ is \textbf{\textit{acyclic}} for the functor $F$ if we have $R^iF(M)=0$ for all $i>0$.
		\end{defn}
		\begin{prop}[4.32]
			Let $M^\bullet$, $i:A\to M^0$ be a resolution of $A$, where the $M^i$ are acyclic for the functor $F$. Then $R^iF(A)$ is equal to the cohomology $H^i(F(M))$ of the complex $F(M^\bullet)$. (\textbf{Obs.} This  cohomology \textit{of a complex} is defined as a cokernel, def. 4.24)
		\end{prop}
		\item (\textbf{Sect. 4.3}) From now on, we consider the category of sheaves of abelian groups over a topological space $X$, and the functor $\Gamma$ of "global sections" which to $\Fc$ associates $\Gamma(X,\Fc)=F(X)$, with values in the category of abelian groups. The category of shaves of abelian groups has sufficiently many injective objects. This implies the existence of the derived functors $R^i\Gamma$. They are written
		\[R^i\Gamma(\Fc)=:H^i(X,\Fc)\]
		\item \begin{defn}[4.33]
			A sheaf $\Fc$ is called \textbf{\textit{flasque}} if for ever pair of open sets $V\subset U$, the restriction map $\rho_{UV}:\Fc(U)\to\Fc(V)$ is surjective.
		\end{defn}
		\begin{prop}[4.34]
			Flasque sheaves are acyclic for the functor $\Gamma$.
		\end{prop}
		\clearpage
		\item \begin{defn}[4.35]
			A \textbf{\textit{fine sheaf}} $\Fc$ over $X$ is a sheaf of $\Ac$-modules, where $\Ac$ is a sheaf of rings over $X$ satisfying the property:
			\begin{quotation}
				For every open cover $U_i$, $i\in I$ of $X$ there exists a partition of unity $f_i$, $i\in I$, $\sum f_i+1$ (where the sum is locally finite), subordinate to this covering.
			\end{quotation}
		\end{defn}
		\begin{prop}[4.36]
			If $\Fc$ is a fine sheaf, we have $H^i(X,\Fc)=0\;\forall i>0$.
		\end{prop}
		\begin{coro}[4.37]
			Let $X$ be a $\Cc^\infty$ manifold. Then
			\[H^k(X,\R)=H^k_{dR}(X,\R).\]
			(The left-hand side is the cohomology of the constant sheaf $\R$, and the right-hand side is the cohomology of the de Rham complex given by exterior differentiation of differential forms)
		\end{coro}
		
		\item \begin{thm}[4.47]
			Let $X$ be a locally contractible topological space. Then we have a canonical isomorphism
			\[H^q_{\sing}(X,\Z)\cong H^q(X,\Z).\]
			The same result holds with $\Z$ replaced by any commutative ring $G$.
		\end{thm}
		\begin{proof}\leavevmode
			\begin{enumerate}
				\item Construct a complex of sheaves from singular cochains. This is done starting with the presheaf
				\[U\mapsto C^q_{\sing}(U,\Z)\]
				where $C^q_{\sing}(U,\Z)$ is just the dual of the $q$-chains, wich in turn are just elements in the free group generated by simplices.
				
				\item Such a complex is a resolution of $\Z$. (computation)
				
				\item Such a resolution is acyclic for $\Gamma$. (prop 4.34; it is flasque)
				
				\item By prop. 4.32 cohomology of acyclic resolution equals derived functors. {\color{blue-violet}Constant sheaf $\Z$ is also acyclic resolution?}
				\[H^q(X,\Z)=H^q(\Gamma(\Cc_{\sing}^\bullet))\]
				
				\item $\Gamma(\Cc^\bullet_{\sing})$ is quasi-isomorphic to $C^\bullet_{\sing}$.
			\end{enumerate}
		\end{proof}\clearpage
		\item \begin{coro}[de Rham theorem]
			\[H_{dR}^q(X)\cong H_{\sing}^q(X,\R)\]
		\end{coro}
		\begin{proof}
			Coro. 4.37 and thm 4.47 replacing $\Z$ with $\R$.
		\end{proof}
		
	\end{enumerate}
	
	\subsection*{Dolbeault's theorem}
	Here's an attempt to use the proof of de Rham's theorem as a model to prove Dolbeault's theorem! 
	\begin{thm}[Dolbeault, coro. 2.6.21 \cite{huybrechts}]
		For a complex compact manifold $M$
		\[H^{p,q}_{\bar\partial}(M)\cong H^q_X(M,\Omega^p)\]
	\end{thm}
	\begin{proof}\leavevmode
		\begin{enumerate}
			\item For every $p$ we have a sheaf complex $\Ac^{p,q}(X)$.
			\item Such a complex is a resolution of the sheaf of sections $\Omega^{p,q}$. (\textbf{Obs.} See prop. 4.19 \cite{voisin}  to use instead the sheaf $\Omega^{0,q}\otimes E$, $E$ holomorphic vector bundle.)
			\item Such a resolution is acyclic. ({\color{blue-violet}maybe by Newlander-Nierenberg+Poincar\'e $\bar\partial$-lemma?)}
			\item By \cite{voisin} prop. 4.32 cohomology of this sheaf complex equals derived functors.
			\item {\color{blue-violet}Dolbeault complex is also a resolution of the sheaf of sections $\Omega^{p,q}$.}
		\end{enumerate}
	\end{proof}
	
\end{document}
